% Part: first-order-logic
% Chapter: models-theories
% Section: expressing-relations

\documentclass[../../../include/open-logic-section]{subfiles}

\begin{document}

\olfileid{fol}{mat}{exr}

\olsection{التعبير عن العلاقات في \article{structure}
  \printtoken{S}{structure}}

\begin{explain}
من أهم استعمالات ال!!{formula}s التعبير عن الخصائص والعلاقات في
!!a{structure}~$\Struct M$ بواسطة أوليات لغة~$\Lang L$ الخاصة
بـ~$\Struct M$. ونعني بهذا ما يأتي: إن !!{domain}~$\Struct M$ مجموعة
من الكائنات. وتؤول ال!!{constant}s وال!!{function}s وال!!{predicate}s
في~$\Struct M$ ببعض الكائنات في~$\Domain M$، والدوال على~$\Domain M$،
والعلاقات على~$\Domain M$، على الترتيب. فإذا كان $\Obj A^2_0$، مثلًا،
في $\Lang L$، أسندت إليه $\Struct M$ علاقة~$R = \Assign{{\Obj
  A^2_0}}{M}$. وعندئذ تعبر الصيغة $\Atom{\Obj A^2_0}{\Obj v_1, \Obj v_2}$
\emph{عن} تلك العلاقة نفسها بالمعنى الآتي: إذا كان إسناد
المتغيرات~$s$ يرسل $\Obj v_1$ إلى $a \in \Domain{M}$ و$\Obj v_2$ إلى
$b \in \Domain M$، فإن
\[
Rab \text{\quad iff\quad} \Sat{M}{\Atom{\Obj A^2_0}{\Obj v_1, \Obj v_2}}[s].
\]
لاحظ أنه لا بد من إدخال إسنادات المتغيرات هنا؛ فلا نستطيع أن نقول
فحسب «$Rab$ إذا وفقط إذا كان $\Sat{M}{\Atom{\Obj A^2_0}{a, b}}$»،
لأن $a$ و$b$ ليسا رمزين من لغتنا، بل هما !!{element}s
في~$\Domain{M}$.

ولأن لدينا أكثر من ال!!{formula}s الذرية، ويمكننا تركيبها باستعمال
الروابط المنطقية والمكممات، تستطيع !!{formula}s أعقد أن تعرّف علاقات
أخرى ليست مضمّنة في~$\Struct M$ مباشرة. ويهمنا أن نعرف كيف نفعل ذلك،
وبوجه خاص أي العلاقات يمكننا تعريفها في !!a{structure}.
\end{explain}

\begin{defn}
لتكن $!A(\Obj v_1,\dots, \Obj v_n)$ !!a{formula} من $\Lang L$ لا
تقع فيها حرة إلا $\Obj v_1$،\dots، $\Obj v_n$، ولتكن $\Struct M$
!!a{structure} لـ~$\Lang L$. تعبر $!A(\Obj v_1,\dots, \Obj v_n)$
\emph{عن العلاقة}~$R \subseteq \Domain M^n$ إذا وفقط إذا كان
\[
Ra_1\dots a_n \text{\quad iff\quad} \Sat{M}{\Atom{!A}{\Obj
    v_1,\dots, \Obj v_n}}[s]
\]
لكل إسناد متغيرات~$s$ يكون فيه $s(\Obj v_i) = a_i$ ($i = 1,
\dots, n$).
\end{defn}

\begin{ex}
في النموذج القياسي للحساب~$\Struct N$، تعبر ال!!{formula} $\Obj
v_1 < \Obj v_2 \lor \eq[\Obj v_1][\Obj v_2]$ عن العلاقة $\le$ على
~$\Nat$. وتعبر ال!!{formula} $\eq[\Obj v_2][\Obj v_1']$ عن علاقة
الخلف، أي العلاقة $R \subseteq \Nat^2$ التي يصدق فيها $Rnm$ إذا كان
$m$ خلف~$n$. وتعبر الصيغة $\eq[\Obj v_1][\Obj v_2']$ عن علاقة السلف.
وتعبر ال!!{formula}s $\lexists[\Obj v_3][(\eq/[\Obj v_3][\Obj 0] \land
  \eq[\Obj v_2][(\Obj v_1 + \Obj v_3)])]$ و$\lexists[\Obj
  v_3][\eq[(\Obj v_1 + {\Obj v_3}')][\Obj v_2]]$ كلتاهما عن العلاقة
$\Obj <$. وهذا يعني أن رمز المحمول~$<$ زائد في الواقع في لغة الحساب؛
إذ يمكن تعريفه.
\end{ex}

\begin{explain}
لا تقتصر أهمية هذه الفكرة على !!{structure}s بعينها، بل تظهر عمومًا
كلما استعملنا لغة لوصف نموذج مقصود أو نماذج مقصودة، أي عندما ننظر في
النظريات. وكثيرًا ما لا تحتوي هذه النظريات إلا عددًا قليلًا من
ال!!{predicate}s بوصفها رموزًا أولية، مع أن علاقات أخرى كثيرة تؤدي
دورًا مهمًا في المجال الذي تستعمل النظريات لوصفه. فإذا أمكن التعبير
عن هذه العلاقات الأخرى على نحو منهجي بواسطة العلاقات التي تؤول
ال!!{predicate}s الأولية للغة، قلنا إننا نستطيع \emph{تعريفها} في
اللغة.
\end{explain}

\begin{prob}
أوجد !!{formula}s في $\Lang L_A$ تعرّف العلاقات الآتية:
\begin{enumerate}
\item $n$ بين $i$ و$j$؛
\item $n$ يقسم $m$ دون باق (أي إن $m$ مضاعف لـ~$n$)؛
\item $n$ عدد أولي (أي لا يقسم~$n$ دون باق عدد غير $1$ و$n$).
\end{enumerate}
\end{prob}

\begin{prob}
افترض أن الصيغة $!A(\Obj v_1, \Obj v_2)$ تعبر عن العلاقة $R
\subseteq \Domain M^2$ في !!a{structure}~$\Struct M$. أوجد صيغًا
تعبر عن العلاقات الآتية:
\begin{enumerate}
\item معكوس العلاقة $R^{-1}$؛
\item الضرب النسبي $R \mid R$؛
\end{enumerate}
هل يمكنك إيجاد طريقة للتعبير عن $R^+$، وهو الإغلاق المتعدي لـ~$R$؟
\end{prob}

\begin{prob}
لتكن $\Lang{L}$ اللغة التي لا تحتوي إلا رمز محمول ثنائي المحل
$<$ (فلا !!{constant}s ولا !!{function}s ولا !!{predicate}s أخرى
فيها---باستثناء~$\eq$ بالطبع). ولتكن $\Struct{N}$ البنية التي فيها
$\Domain{N} = \Nat$ و$\Assign{<}{N} = \Setabs{\tuple{n,m}}{n <
  m}$. أثبت ما يأتي:
\begin{enumerate}
\item $\{ 0 \}$ قابلة للتعريف في $\Struct{N}$؛
\item $\{ 1 \}$ قابلة للتعريف في $\Struct{N}$؛
\item $\{ 2 \}$ قابلة للتعريف في $\Struct{N}$؛
\item لكل $n \in \Nat$، المجموعة $\{ n \}$ قابلة للتعريف في
  $\Struct{N}$؛
\item كل مجموعة جزئية منتهية من $\Domain{N}$ قابلة للتعريف في
  $\Struct{N}$؛
\item كل مجموعة جزئية من $\Domain{N}$ متممتها منتهية قابلة للتعريف
  في $\Struct{N}$ (حيث تكون متممة $X \subseteq \Nat$ منتهية إذا
  وفقط إذا كانت $\Nat \setminus X$ منتهية).
\end{enumerate}
\end{prob}


\end{document}
